\documentclass[border=0pt]{standalone}
% The shared preamble of every docs diagram (docs/diagrams/<name>.tex).
% dvisvgm's driver first, so patterns and opacity survive into the SVG.
\def\pgfsysdriver{pgfsys-dvisvgm.def}
\usepackage[siunitx,RPvoltages]{circuitikz}
\usepackage{xcolor}
% Latin Modern: outline fonts for every glyph a label uses, µ included.
\usepackage[T1]{fontenc}
\usepackage{lmodern}
\usetikzlibrary{arrows.meta,backgrounds,calc,patterns,positioning}
\tikzset{
  % A white card, so diagrams read in mdBook's dark themes too.
  every picture/.append style={show background rectangle,
    background rectangle/.style={fill=white,inner sep=4pt}},
  % Layout geometry: layers told apart by fill and hatching, never colour alone.
  layer a/.style={draw=black,line width=0.6pt,fill=gray!20},
  layer b/.style={draw=black,line width=0.6pt,fill=gray!45},
  layer c/.style={draw=black,line width=0.6pt,fill=white,postaction={pattern=north east lines}},
  layer d/.style={draw=black,line width=0.6pt,fill=white,postaction={pattern=crosshatch dots}},
  marker/.style={draw=black,dashed,line width=0.6pt},
  % A measured distance and its label.
  dim/.style={{Latex[length=2mm]}-{Latex[length=2mm]},line width=0.5pt},
  dim label/.style={font=\small,fill=white,inner sep=1pt},
  % Where a rule fails: thick, dashed and labelled, not only red.
  violation/.style={draw=red!75!black,line width=1.4pt,dash pattern=on 3pt off 2pt},
  violation label/.style={font=\small\bfseries,text=red!75!black},
  caption/.style={font=\small},
}

\begin{document}
\begin{circuitikz}[font=\small,
    dev/.style={draw,rectangle,minimum width=1.5cm,minimum height=0.7cm,fill=gray!20},
    net/.style={draw,circle,minimum size=0.85cm,inner sep=1pt},
    role/.style={font=\footnotesize,fill=white,inner sep=1pt}]
  % Left: the inverter as drawn in a schematic.
  \node[font=\bfseries] at (1.2,5.2) {Inverter};
  \node[pmos] (p) at (1.2,3.0) {};
  \node[right=2pt] at (p.east) {M2 pfet};
  \draw (p.S) -- ++(0,0.3) node[vcc] {VDD};
  \node[nmos] (n) at (1.2,0.9) {};
  \node[right=2pt] at (n.east) {M1 nfet};
  \draw (p.D) -- (n.D);
  \draw (p.D) ++(0,-0.3) coordinate (y) -- ++(1.3,0) node[right] {Y};
  \draw (n.S) node[ground] {};
  \node[right] at ($(n.S)+(0.25,-0.3)$) {VSS};
  \draw (p.G) -- ++(-0.5,0) coordinate (a) -- (a |- n.G) -- (n.G);
  \draw (a |- p.G) node[left] {A};

  % Right: the same circuit as GPurify's graph.
  \node[font=\bfseries] at (7.6,5.2) {As a graph};
  \node[net] (vdd) at (7.6,3.7) {VDD};
  \node[dev] (m2) at (7.6,2.5) {M2 pfet};
  \node[net] (a2) at (5.4,1.6) {A};
  \node[net] (y2) at (9.8,1.6) {Y};
  \node[dev] (m1) at (7.6,0.7) {M1 nfet};
  \node[net] (vss) at (7.6,-0.5) {VSS};
  \draw (m2) -- node[role] {channel} (vdd);
  \draw (m2) -- node[role,pos=0.45] {gate} (a2);
  \draw (m2) -- node[role,pos=0.45] {channel} (y2);
  \draw (m1) -- node[role,pos=0.45] {gate} (a2);
  \draw (m1) -- node[role,pos=0.45] {channel} (y2);
  \draw (m1) -- node[role] {channel} (vss);
  \node[caption,align=left] at (7.6,-1.4) {devices are boxes, nets are circles,\\ each edge a terminal with its role};
\end{circuitikz}
\end{document}
